% 01-introduction.tex
\section{Introduction}
\label{sec:intro}

TensorPlay is a readability-first implementation of a dense tensor core, in approximately 75\,k lines of C++ (plus generated bindings) and several hundred Python modules under
. The build produces four shared libraries \texttt{p10}, \texttt{tpx}, \texttt{stax}, \texttt{tp\_python} and one extension \texttt{\_C}
 whose \texttt{OUTPUT\_DIRECTORY} is pinned by to \texttt{TP\_PACKAGE\_DIR=tensorplay}
. Throughout the paper, every mechanism is cited as \texttt{path:line} so that the text and the source stay one jump apart.

\subsection{Problem: Readability vs.\ Performance}
\label{sec:intro:problem}

Mainstream frameworks co-locate tensor storage, operator dispatch, automatic differentiation (\ac{ad}), compilation, and distribution in one translation unit. The result is maximal visibility for
the optimizer: a single \texttt{LTO} step can inline storage access into the kernel loop. The cost is cognitive and incremental-build overhead. Understanding how a tensor is stored
requires traversing an inheritance chain (\texttt{Tensor} $\rightarrow$ \texttt{Variable} $\rightarrow$ storage) and the associated intrusive reference counting; modifying a CPU kernel
for \texttt{add} triggers a rebuild of the entire operator aggregation target; implicit header inclusion creates de~facto circular dependencies between core and operator layers
that static analysis cannot recover without building the graph first.

TensorPlay inverts the trade-off explicitly. It optimizes for \emph{readability}: a newcomer can understand tensor storage by reading only \texttt{p10}
 without loading any differentiation concept; an optimizer for
\texttt{add} rebuilds only the tensor library (one build target, one relink); a compiler experiment lives inside \texttt{stax}
without relinking storage.
The price is an explicit cross-library boundary per call. Concretely, one \texttt{atomic<void*>} acquire load, one virtual dispatch, and no \ac{lto} whole-program inlining across the four pillars.
In short, calls must serialize across the boundary, but the gains in readability and incremental build are larger than the cost.
The project accepts this cost because the target workloads are teaching, correctness inspection, and hardware experimentation---not peak FLOPS on a frozen kernel set.

Concretely, the boundary is enforced at link time, not by convention. The base library declares a pure interface (\texttt{requires\_grad}, \texttt{grad}, \texttt{set\_grad}); the tensor implementation holds \texttt{shared\_ptr<AutogradMetaBase> autograd\_meta\_} that is deliberately omitted from the copy constructor, so \texttt{Tensor b=a} shares storage, sizes, strides, device, dtype and version counter but not history. The autograd library supplies the concrete plug (\texttt{grad\_fn\_}, \texttt{grad\_accumulator\_} as \texttt{weak\_ptr<Node>} to avoid cycles, view replay state); one boundary file performs the single \texttt{static\_cast<AutogradMeta*>} handshake. Because no header in the base pillar sees the autograd pillar, adding that include inside the tensor pillar fails with \texttt{undefined reference} the link step.

Build granularity follows the same axiom. The build script adds the four pillars in dependency order---tensor core, autograd, compiler---and each links downward only. Touching a CPU kernel file dirties exactly \texttt{libp10.so} and the final \texttt{\_C} link; touching the engine file rebuilds \texttt{libtpx.so} plus \texttt{libstax.so} and \texttt{\_C} but not \texttt{libp10.so}. Each library locates its siblings through an \texttt{RPATH} of \texttt{\$ORIGIN} (the extension uses \texttt{\$ORIGIN/../lib}). On Windows the same invariant is preserved by an ordered \texttt{LoadLibraryExW} preload walk.

\subsection{Contributions}
\label{sec:intro:contributions}

This paper makes four concrete contributions beyond high-level architecture sketches. Each is a mechanism with a working implementation, not a design note.
Table~\ref{tab:intro:contrib} is the index; \S\ref{sec:intro:c1}--\S\ref{sec:intro:c4} expand each row.

\begin{table}[H]
\centering
\caption{Four concrete contributions. Every entry cites the line that implements the mechanism, not a design note.}
\label{tab:intro:contrib}
\small
\begin{tabular}{p{0.6cm}p{2.6cm}p{4.0cm}p{4.6cm}}
\toprule
\# & Mechanism & Primary \texttt{path:line} & Paper section \\
\midrule
C1 & Four-pillar acyclic graph &, &
\S\ref{sec:arch} \\
C2 & 13-slot redispatch + registration &,
\texttt{CPythonBridge.cpp} & \S\ref{sec:redispatch} \\
C3 & Explicit DAG without Tape & \texttt{Node.h}, \texttt{Engine.h}, \texttt{SavedVariable.h}, &
\S\ref{sec:autograd} \\
C4 & Vmap at offset~9 + profiler stack & \texttt{DispatchKey.h}, \texttt{Profiler.h} &
\S\ref{sec:vmap}, \S\ref{sec:profiler} \\
\bottomrule
\end{tabular}
\end{table}

\subsubsection{C1: Four-Pillar Decomposition with Link-Time Acyclicity}
\label{sec:intro:c1}

The framework is partitioned into \texttt{p10} (tensor, storage, 300+ kernels, allocators), \texttt{tpx} (explicit DAG, engine, saved variables), \texttt{stax} (minimal \texttt{Graph} IR with a single \texttt{fused\_mul\_add} pass), and the Python extension \texttt{\_C} plus its single-copy bridge library \texttt{tp\_python}. Table~\ref{tab:intro:pillars} repeats the directory $\rightarrow$ artifact contract.
\begin{table}[H]
\centering
\caption{Directory $\rightarrow$ CMake target $\rightarrow$ artifact, with link dependencies (downward only).}
\label{tab:intro:pillars}
\small
\begin{tabular}{p{2.0cm}p{1.4cm}p{3.6cm}p{4.8cm}}
\toprule
Directory & Target & Artifact & Link deps \\
\midrule
core & \texttt{p10} & \texttt{libp10.so} & \texttt{Threads}, \texttt{COMPUTE\_LIBS} \\
autograd & \texttt{tpx} & \texttt{libtpx.so} & \texttt{p10} \\
compiler & \texttt{stax} & \texttt{libstax.so} & \texttt{p10}, \texttt{tpx} \\
bridge & \texttt{tp\_python} & \texttt{libtp\_python.so} & \texttt{p10}, \texttt{pybind11} \\
extension & \texttt{\_C} & \texttt{\_\_init\_\_.so} & \texttt{p10}, \texttt{tpx}, \texttt{stax}, \texttt{tp\_python} \\
\bottomrule
\end{tabular}
\end{table}

Axiom: $\forall (a,b)\in D$, \texttt{level}(a)$>$\texttt{level}(b) with \texttt{level}(p10)=0, \texttt{level}(tpx)=1, \texttt{level}(stax)=2, \texttt{level}(\_C)=3 and
\texttt{level}(tp\_python)=3. The enforcements are (i) directory isolation, (ii) \texttt{target\_link\_libraries} edges that point only downward, and (iii) header visibility: no
\texttt{*.h} includes \texttt{tpx}. The recorded move at records that concrete
metadata once lived in \texttt{p10} and was vacated to preserve the axiom. The Windows static branch at embeds the pillars into \texttt{\_C.pyd} to
avoid \texttt{\_\_declspec(dllexport)} on every p10 symbol while keeping visibility identical.

\subsubsection{C2: Redispatch and Static Registration}
\label{sec:intro:c2}

 encodes 13 keys: \texttt{CPU}=0, \texttt{CUDA}=1, \texttt{Vulkan}=2, \texttt{AutogradCPU}=3--5, \texttt{AutocastCPU}=6--8,
\texttt{VmapCPU}=9--11, \texttt{Composite}=12, \texttt{EndOfKeys}=13. Offsets \texttt{kAutogradKeyOffset}=3, \texttt{kAutocastKeyOffset}=6, \texttt{kVmapKeyOffset}=9
 make priority arithmetic. \texttt{DispatchKeySet} is a \texttt{uint32\_t} bitset whose \texttt{highest\_priority\_key}
at scans the most-significant set bit so numerically larger keys win; \texttt{TensorImpl::key\_set} at computes the
set from device and \texttt{inference\_tensor\_} without storing it.

 stores \texttt{array<atomic<KernelFunction>,13>} per operator; \texttt{OperatorHandle::getKernel} at does one
\texttt{acquire} load and, only for \texttt{is\_backend\_key} (0--2), a second load at \texttt{Composite} index~12. \texttt{DispatchStub::call} is the
single choke point that enforces the autocast contract \texttt{is\_autocast\_key(k) \&\& !is\_enabled(k) $\rightarrow$ throw} with \texttt{[[unlikely]]} and emits a 32-frame
\texttt{backtrace} when \texttt{TP\_TRACE\_KERNEL\_MISS} is set. Registration is bulk: \texttt{TENSORPLAY\_LIBRARY\_IMPL(KEY, NAME)} uses two-level
\texttt{TP\_CONCAT} at to produce hygienic unique statics; \texttt{Library::impl} at chains calls. The deprecated per-op
\texttt{TENSORPLAY\_REGISTER\_KERNEL} is retained as a legacy registration path.

Python binding is dual-path and fill-only. The bridge and the generated module together implement parse (
\texttt{tpx\_py\_parse\_into} via \texttt{memcpy} + linear keyword scan), check (16 kinds + \texttt{TPK\_OPTIONAL}=0x80), subclass/function hooks
(via \texttt{PythonDispatchTLS}), unpack (fast \texttt{instance} holder), and wrap (
\texttt{g\_wrap\_cache} with capsule destructor). The codegen emits \texttt{generated\_functions[]} and
\texttt{register\_generated\_cpython\_functions} that never shadows a \texttt{pybind11} name (checked by \texttt{PyObject\_HasAttrString}).
It installs the generated table last behind \texttt{TP\_NO\_FASTCALL}; the bridge keeps a single \texttt{tp\_python} so the caster registry and wrap cache are not duplicated
per extension.

\subsubsection{C3: Autograd}
\label{sec:intro:c3}

There is no \ac{tape}. Each forward op that requires grad allocates a \texttt{Node} subclass (148 types generated from via into \texttt{AutogradNodesGenerated.h}, or a handwritten node in).
\texttt{next\_edges: vector<Edge>} at points to parents; \texttt{Edge} at carries a \texttt{weak\_ptr<Node>} plus \texttt{output\_nr}.
\texttt{output\_nr}. The graph is the transitive closure of those edges. \texttt{Node::sequence\_nr} at \texttt{Node.h} is a global monotonic counter from
\texttt{get\_and\_increment\_sequence\_nr}, used by the engine to prioritize recomputation order deterministically.

\texttt{GraphTask.h} holds the task launched by \texttt{backward}: \texttt{dependencies\_} in-degree counts, \texttt{captured\_vars}, and the initiating \texttt{GraphRoot}
synthesized at \texttt{Engine.cpp} only when \texttt{roots.size!=1}. is the per-node accumulator whose \texttt{accumulate} adds
incoming grads in topological order. defines \texttt{ReadyQueue} as a thread-safe max-heap by \texttt{sequence\_nr} ( \texttt{return
fn->sequence\_nr < other.fn->sequence\_nr}); dispatches with device-partitioned queues (\texttt{ready\_queues\_},
\texttt{queue\_for\_device}, \texttt{execute\_task} at) and a \texttt{local\_queue} for reentrancy so that \texttt{Node} implementations
that call \texttt{backward} recursively do not deadlock the global worker.

Versioning is fail-fast, not silent. \texttt{VariableVersion.h} backs \texttt{TensorImpl::version\_counter\_} at with \texttt{bump}
 on every inplace mutation. \texttt{SavedVariable.h} saves \texttt{saved\_version\_}; the unpack checks \texttt{expected==current} and throws
rather than producing a numerically wrong gradient, which is the single most valuable readability guard in the engine. additionally stores
\texttt{grad\_accumulator\_} as \texttt{weak\_ptr<Node>} so leaf accumulators do not form uncollectable cycles.

\begin{tpinsight}[Why no Tape?]
A Wengert-list tape records every primal operation into one contiguous log and replays it in reverse. TensorPlay instead materializes a \emph{node graph}: each differentiable op allocates its own \texttt{Node} whose \texttt{next\_edges} already point to parents, so the ``replay'' is a topological walk over pointers that already exist. The tape variant would centralize bookkeeping in one object---attractive for a JIT that wants to fuse the whole log---but it forces every op to interact with a global recorder, breaking the pillar axiom (\S\ref{sec:intro:c1}) because \texttt{p10} kernels would need to know about the recorder's memory. The node-graph variant pays one allocation per op instead, which \S\ref{sec:eval} shows is invisible on any workload where kernels dominate, and it keeps the recorder logic entirely inside \texttt{tpx}.
\end{tpinsight}

\begin{tpdef}[Definition (Link-Time Acyclicity)]
Let $V$ be the set of build targets and let $D \subseteq V \times V$ be the ``links'' relation induced by the build system's \texttt{target\_link\_libraries} directives. The build is \emph{acyclic} iff there exists a level function $\mathrm{level}:V\rightarrow\mathbb{N}$ with $\forall(a,b)\in D:\ \mathrm{level}(a)>\mathrm{level}(b)$. TensorPlay fixes \texttt{level(p10)=0}, \texttt{level(tpx)=1}, \texttt{level(stax)=2}, \texttt{level(\_C)=3}; the property is enforced by build-target ordering and certified by the vacated-symbol move.
\end{tpdef}

\subsubsection{C4: Vmap and the Profiler Stack}
\label{sec:intro:c4}

Vmap occupies offset~9 precisely so it outranks autograd and backend. \texttt{DispatchKey::VmapCPU}=9 at with \texttt{kVmapKeyOffset}=9
; \texttt{toVmapKey} at and \texttt{is\_vmap\_key} are \texttt{constexpr} range tests that lower to one
compare. \texttt{TransformDispatch} at and the kernel table at implement per-slice fallback versus
batched kernels; the dispatcher unwraps transform state before the autograd layer observes the tensor, so \texttt{Vmap} and autograd compose without a second dispatch table.

The profiler spans \ac{cpu} and \ac{cuda} with one hot-path guard. exports \texttt{g\_active}, \texttt{g\_capture\_shapes}, \texttt{g\_capture\_sites} as
\texttt{atomic<bool>}; when inactive the cost per op is one acquire load. \texttt{OpRecord} at stamps a slot with \texttt{end\_ns}, \texttt{tid},
\texttt{kind}, optional \texttt{dtypes} and \texttt{stack\_id}; \texttt{GpuTimerPair} at arms a pooled \texttt{cudaEvent} pair per dispatched
CUDA op and resolves elapsed \texttt{gpu\_ms} at session stop. \ac{cupti} \texttt{stop} at \texttt{Profiler.h} dlopen \texttt{libcupti} and correlate kernels to
ops via external correlation ids pushed in \texttt{GpuTimerPair}; \ac{nvtx} \texttt{pop} at and \ac{itt} \texttt{itt\_task\_begin\_name}
 bridge to Nsight and VTune. All collectors degrade gracefully when the native library is absent, and the binding layer at \texttt{profiler.py}
multiplexes sampling without adding a build dependency.

\subsection{Scope and Non-goals}
\label{sec:intro:scope}

Table~\ref{tab:intro:scope} documents what is implemented and deliberately leave other axes as \texttt{Composite} fallbacks or as not implemented. Table~\ref{tab:intro:scope} enumerates
the trimmed axes; every absent row is intentional for the 75\,k-line readability target.

\begin{table}[H]
\centering
\caption{Scope: present versus trimmed. Trimmed means no dispatch key and no dedicated storage; the generic path remains valid via \texttt{Composite}.}
\label{tab:intro:scope}
\small
\begin{tabular}{p{2.8cm}p{4.6cm}p{4.6cm}}
\toprule
Axis & Present in & Trimmed (non-goal) and fallback \\
\midrule
Dispatch keys & 13 keys at (\texttt{Vulkan} + \texttt{Vmap} + \texttt{Composite}) & No sparse, quantized, meta, mkldnn, xla, lazy
dispatch keys; quantized kernels exist only as composites \\
Tensor storage & Dense \texttt{Storage}, \texttt{DataPtr}, \texttt{Allocator}
 & No sparse COO/CSR storage, no quantized per-tensor/per-channel storage, no meta-device allocation; sparse helpers
are teaching stubs \\
Autograd & Explicit DAG (, \texttt{GraphTask.h},) with \texttt{SavedVariable} version guard & No
tape reuse, no boxed alias analysis, no inference-mode fast path beyond \texttt{InferenceMode} at \\
Compiler & Minimal \texttt{Graph} IR at with vertical \texttt{fused\_mul\_add} & No loop fusion, layout planner, vectorization
policy, auto-tuning, or persistence beyond in-memory \texttt{Graph::execute} at \\
Distribution & Single-process collectives via \texttt{tensorpipe} link at & No sharded \texttt{DTensor}, no pipeline parallelism, no heterogeneous device
mesh \\
Quantization & \texttt{Quantizer} at placeholder & No quantized dispatch axis, no kernel selection by \texttt{QScheme}; 8-bit kernels are composite
demos \\
\bottomrule
\end{tabular}
\end{table}

Sparse is the most illustrative trim. declares the API surface, but never defines \texttt{SparseCPU} and
\texttt{TensorImpl::key\_set} never emits it, so dispatch never queries a sparse slot. A contributor adding sparse would add a key, extend \texttt{kDispatchKeyCount} from 13, and add
the offset helpers, which is exactly the change the acyclic build is designed to make painful enough to discuss. Meta tensors are absent for the same reason: \texttt{DeviceType::Meta}
exists as an enum value but never participates in \texttt{computeDispatchKey} at and never acquires a \texttt{RegisterKernel} entry, so \texttt{has\_kernel(...,
Meta)} remains false by construction. Boxed alias analysis and shape-check elision are absent because the codegen pipeline at and
 deliberately generates straight-line backward formulas without alias metadata; the stax interpreter re-enters \texttt{tpx::ops}
 so correctness of aliasing stays with the eager dispatcher rather than with a compiler pass.

The trimmed axes are not undocumented; they are documented as trimmed so that a reader does not search for a sparse dispatch key that does not exist. A future revision that promotes a trimmed axis to a real key must update and.

\subsection{How to Read This Paper}
\label{sec:intro:howto}

Figure~\ref{fig:pillars} is the map of the whole white paper; every later figure is a zoom of one arrow in it. The paper supports both a linear reading and a topic-first reading; Table~\ref{tab:intro:guide} is the topic-first index.

\begin{table}[H]
\centering
\caption{Reader's guide. Each row names the chapters and central figure for one learning goal.}
\label{tab:intro:guide}
\small
\begin{tabular}{p{3.2cm}p{4.2cm}p{4.2cm}}
\toprule
Goal & Chapters & Central figure \\
\midrule
Understand the pillars & \S\ref{sec:arch} & Fig.~\ref{fig:pillars} \\
Follow one op end-to-end & \S\ref{sec:redispatch}, \S\ref{sec:memory} & Figs.~\ref{fig:dispatchermap}, \ref{fig:e2e} \\
Understand backward & \S\ref{sec:autograd} & Fig.~\ref{fig:usb}, the DAG at \\
Understand batched AD & \S\ref{sec:vmap} & The vmap offset diagram at \\
Follow code generation & \S\ref{sec:codegen} & Fig.~\ref{fig:codegen_contract} \\
Read the cost of choices & \S\ref{sec:eval} & Tables and histograms there \\
\bottomrule
\end{tabular}
\end{table}

System overview: read \S\ref{sec:arch} and Fig.~\ref{fig:pillars} first; the YAML$\rightarrow$artifact contract underneath it is at via \texttt{main.py}. Dispatch internals live in
\S\ref{sec:redispatch} (DispatchKey arithmetic, \texttt{array<atomic<void*>,13>}, Composite fallback, \texttt{TP\_CONCAT}) and \S\ref{sec:memory} (TensorIterator seven-step pipeline
). Autograd is the subject of \S\ref{sec:autograd}: the USB slot at $\rightarrow$
$\rightarrow$ with cast at \texttt{Autograd.cpp}, the heap \texttt{ReadyQueue} by \texttt{sequence\_nr}, and the
\texttt{SavedVariable} guard. The compiler pillar is \S\ref{sec:arch:stax} and the stax IR at plus the interpreter; every other Python
module ultimately re-enters \texttt{\_C}, so symbol renaming and lazy \texttt{\_\_getattr\_\_}
are the only Python-side complexity with cross-cutting impact.

\begin{figure}[H]
\centering
\begin{tikzpicture}[node distance=0.6cm]
\node[pillar, fill=blue!8, draw=tpblue] (p10) {p10\\{\footnotesize Tensor / Storage}\\{\footnotesize Dispatcher / Kernels}};
\node[pillar, fill=orange!8, draw=tporange, right=0.7cm of p10] (tpx) {tpx\\{\footnotesize Autograd DAG}\\{\footnotesize Engine / Node}};
\node[pillar, fill=purple!8, draw=tppurple, right=0.7cm of tpx] (stax) {stax\\{\footnotesize Graph IR}\\{\footnotesize Pass / Interpreter}};
\node[pillar, fill=green!8, draw=tpgreen, right=0.7cm of stax] (c) {\_C + tp\_python\\{\footnotesize pybind11 + FASTCALL}};
\node[libbox, fill=blue!4, draw=tpblue, below=0.7cm of p10, minimum width=13.2cm, minimum height=0.7cm] (base) {CMake + codegen \quad \texttt{native\_functions.yaml $\rightarrow$ 10 artifacts} \quad (schema parser)};
\node[libbox, fill=yellow!12, draw=orange!60!black, above=0.7cm of p10, minimum width=13.2cm, minimum height=0.7cm] (py) {Python veneer \quad \texttt{ graph}};
\draw[arrow, color=gray] (p10) -- (tpx);
\draw[arrow, color=gray] (tpx) -- (stax);
\draw[arrow, color=gray] (stax) -- (c);
\foreach \x in {p10,tpx,stax,c} \draw[arrow, dashed, color=orange!70!black] (\x) -- (py);
\foreach \x in {p10,tpx,stax,c} \draw[arrow, dashed, color=tpblue!70] (base) -- (\x);
\end{tikzpicture}
\caption{Four pillars and dependency direction (acyclic). Downward edges only; any upward include fails at link time. Solid arrows: link dependency; dashed: generation or re-export.}
\label{fig:pillars}
\end{figure}

\subsection{Conventions}
\label{sec:intro:conventions}

\texttt{path:line} means the file relative to the repository root; is line~49 of that file, and \texttt{path:line--line} is an
inclusive range. Code listings are verbatim excerpts with at most whitespace folding; the caption gives the exact range. Tables enumerate normative values (key numbers, artifact names, trimmed axes) rather than
illustrative samples. \texttt{Composite} is a backend-neutral fallback; it never appears in a \texttt{DispatchKeySet} and is consulted only by
\texttt{OperatorHandle::getKernel} for \texttt{is\_backend\_key}. \texttt{TP\_CONCAT} is a two-level macro so its argument expands before pasting;
the paper always cites the two-level definition, not a single-level variant.

\begin{tppitfall}[Pitfall: Searching for a Dispatch Key That Is Never in the Set]
A reader tracing dispatch logic often searches for \texttt{Composite} inside \texttt{DispatchKeySet} construction and finds nothing---and concludes composite fallback is dead code. It is not: \texttt{Composite} is \emph{never} a member of any key set. It is consulted only as a fallback inside \texttt{OperatorHandle::getKernel} when \texttt{is\_backend\_key(key)} holds and the backend slot is \texttt{nullptr}. The same trap recurs for \texttt{Meta}: the \texttt{DeviceType} enum carries the value, but \texttt{computeDispatchKey} never emits it. When tracing dispatch, always start from the key set the tensor actually carries, not from the key enum.
\end{tppitfall}

\subsection{Design Trade-offs}
\label{sec:intro:tradeoffs}

The readability-first stance rejects three alternatives that a production framework would accept.

\textbf{Monolithic single library.} The obvious alternative to four pillars is one \texttt{libtensorplay.so} with unity-build speed and cross-module inlining. It was rejected because every kernel edit would relink the world: the per-library granularity in \S\ref{sec:intro:c1} (touch one CPU kernel, relink only \texttt{libp10.so}) is the primary developer-experience claim of this project, and it structurally requires separate link units. The measured cost of the split---one atomic load and one virtual call per op, quantified in \S\ref{sec:eval}---is a constant $\sim$2\,ns that any real kernel dwarfs.

\textbf{Intrusive reference counting.} Storage and graph nodes could use intrusive pointers to eliminate the control-block allocation and halve cache traffic on \texttt{shared\_ptr} copies. Rejected for observability: intrusive counts cannot be inspected in \texttt{gdb} without pretty-printers, and the ownership question ``who holds this storage?'' is exactly what a teaching codebase must answer with standard tools. \S\ref{sec:eval} shows the constant-factor cost is invisible at kernel scale.

\textbf{Runtime extensibility.} Third-party dispatch keys registered at import time would make the dispatcher a plugin platform. Rejected because it converts the 13-slot array into a sparse map, raising every \texttt{getKernel} from an array index to a hash lookup, and because it moves op registration from reviewable YAML into unreviewable \texttt{\_\_init\_\_} side effects. The closed-key design keeps registration inspectable in one file.

